% ===== PRÉAMBULE CANONIQUE — à coller VERBATIM en tête de CHAQUE .tex =====
\documentclass[11pt,a4paper]{article}
\usepackage[utf8]{inputenc}\usepackage[T1]{fontenc}\usepackage{lmodern}
\usepackage[french]{babel}
\usepackage{amsmath,amssymb,mathtools}\usepackage{icomma}
\usepackage{siunitx}\sisetup{locale=FR}  % \qty{}{}, \si{}, \ang{}
\usepackage{xcolor}\usepackage{colortbl}
\usepackage{tikz}\usepackage{pgfplots}\pgfplotsset{compat=1.18}
\usepackage[siunitx,europeanresistors]{circuitikz} % circuits (élec.)
\usepackage[most]{tcolorbox}\usepackage{enumitem}\usepackage{array,booktabs}
\usepackage[a4paper,margin=2.2cm]{geometry}
\usetikzlibrary{arrows.meta,calc,angles,quotes,positioning,patterns,%
  backgrounds,shapes.geometric,decorations.pathmorphing,decorations.markings,intersections}
\definecolor{primary}{RGB}{222,49,99}\definecolor{primarydark}{RGB}{150,22,55}
\definecolor{defcol}{RGB}{0,114,178}\definecolor{methcol}{RGB}{52,92,148}
\definecolor{excol}{RGB}{0,135,90}\definecolor{attncol}{RGB}{200,40,55}
\tcbset{boxbase/.style={enhanced,breakable,boxrule=0.9pt,arc=2.5pt,
  left=3mm,right=3mm,top=2.3mm,bottom=2.3mm,fonttitle=\bfseries,coltitle=white}}
% --- boîtes TUTORIEL (noms EXACTS, ne pas renommer) ---
\newtcolorbox{cle}[1][]{boxbase,colback=defcol!6,colframe=defcol,
  title={\textsc{À retenir}\ifx\relax#1\relax\relax\else\textmd{ : \itshape #1}\fi}}
\newtcolorbox{methode}[1][]{boxbase,colback=methcol!7,colframe=methcol,sharp corners,
  title={\textsc{Méthode}\ifx\relax#1\relax\relax\else\textmd{ --- \itshape #1}\fi}}
\newtcolorbox{exemplebox}[1][]{boxbase,colback=excol!5,colframe=excol,
  title={\textsc{Exemple}\ifx\relax#1\relax\relax\else\textmd{ : \itshape #1}\fi}}
\newtcolorbox{piege}[1][]{boxbase,colback=attncol!6,colframe=attncol,
  title={\textsc{Piège}\ifx\relax#1\relax\relax\else\textmd{ : \itshape #1}\fi}}
\newtcolorbox{remarque}[1][]{boxbase,colback=black!4,colframe=black!45,
  title={\textsc{Remarque}\ifx\relax#1\relax\relax\else\textmd{ : \itshape #1}\fi}}
% --- boîtes EXERCICE / CORRIGÉ (noms EXACTS, auto-comptées) ---
\newtcolorbox[auto counter]{exo}[1][]{boxbase,colback=primary!4,colframe=primary,
  title={\textsc{Exercice}~\thetcbcounter\ifx\relax#1\relax\relax\else\textmd{ --- \itshape #1}\fi}}
\newtcolorbox[auto counter]{corrige}[1][]{boxbase,colback=excol!4,colframe=excol,
  title={\textsc{Corrigé}~\thetcbcounter\ifx\relax#1\relax\relax\else\textmd{ --- \itshape #1}\fi}}
\newcommand{\vv}[1]{\vec{#1}}\newcommand{\ds}{\displaystyle}
\newcommand{\R}{\mathbb{R}}\newcommand{\N}{\mathbb{N}}\newcommand{\Z}{\mathbb{Z}}
% (un \usepackage{fancyhdr}+en-tête est possible ; il est ignoré côté web)
% =========================================================================
\begin{document}
\begin{corrige}[barycentre de deux masses]
\begin{enumerate}[label=\alph*)]
  \item $\dfrac{\overline{AG}}{\overline{BG}}=\dfrac{m_2}{m_1}=\dfrac{6}{2}=3$, donc
        $\overline{AG}=3\,\overline{BG}$. Avec $\overline{AG}+\overline{BG}=80$ :
        $4\,\overline{BG}=80 \Rightarrow \overline{BG}=\qty{20}{cm}$, $\overline{AG}=\qty{60}{cm}$.
        ($G$ est donc à \qty{60}{cm} de $A$.)
  \item $G$ est à \qty{20}{cm} de $m_2$ contre \qty{60}{cm} de $m_1$ : il est bien plus proche de la
        masse la plus lourde $m_2$.
\end{enumerate}
\end{corrige}

\begin{corrige}[trois masses alignées]
\begin{enumerate}[label=\alph*)]
  \item $x_{G_{12}} = \dfrac{m_1 x_1 + m_2 x_2}{m_1+m_2} = \dfrac{1\times 0 + 2\times 30}{3}
        = \qty{20}{cm}$ ; ce point porte $m_1+m_2=\qty{3}{kg}$.
  \item On compose $G_{12}$ (\qty{3}{kg} en $x=\qty{20}{cm}$) avec $m_3$ (\qty{3}{kg} en
        $x=\qty{60}{cm}$) :
        \[ x_G = \frac{3\times 20 + 3\times 60}{3+3} = \frac{240}{6} = \qty{40}{cm}. \]
        (On retrouve bien le calcul direct $\tfrac{1\cdot 0+2\cdot 30+3\cdot 60}{6}=\qty{40}{cm}$.)
\end{enumerate}
\end{corrige}

\begin{corrige}[une tige lestée]
\begin{enumerate}[label=\alph*)]
  \item On traite la tige comme une masse $M$ concentrée en son centre ($x=\qty{50}{cm}$) :
        \[ x_G = \frac{M\times 50 + m\times 100}{M+m} = \frac{2\times 50 + 3\times 100}{2+3}
        = \frac{400}{5} = \qty{80}{cm}. \]
        Le centre de gravité se déplace vers la masse ajoutée.
\end{enumerate}
\end{corrige}

\begin{corrige}[méthode des deux suspensions]
\begin{enumerate}[label=\alph*)]
  \item À l'équilibre, $G$ se place \emph{à la verticale} du point de suspension : il est donc
        \textbf{sur la verticale tracée depuis $P_1$} (quelque part le long de cette droite).
  \item On recommence en suspendant la plaque par un \emph{autre} point $P_2$ et on trace la
        nouvelle verticale : $G$ est à l'\textbf{intersection} des deux droites. \textbf{Deux}
        suspensions suffisent.
\end{enumerate}
\end{corrige}

\begin{corrige}[stabilité : large ou haut ?]
\begin{enumerate}[label=\alph*)]
  \item L'objet ne bascule pas tant que la \textbf{verticale passant par $G$} tombe \emph{à l'intérieur
        de sa base d'appui}.
  \item L'objet \textbf{A} (large et bas) est le plus stable : son $G$ est bas et sa base large, donc il
        faut l'incliner \emph{beaucoup} avant que la verticale de $G$ ne sorte de la base. L'objet B
        (étroit et haut) a un $G$ élevé au-dessus d'une base étroite : une faible inclinaison suffit à
        faire sortir la verticale de $G$ de la base, et il bascule.
\end{enumerate}
\end{corrige}

\end{document}
